\subsection{THE CASE OF A LINEAR FAMILY OF ENDOMORPHISMS}

Assume now that S has a vector space structure, that the map \(r : S \to \text{End}(V)\) is linear, and that V and S are finite dimensional.

PROPOSITION 5. Assume that condition (AC) is satisfied, and let \(\lambda : S \to k\) be such that \(V^{\lambda}(S) \neq 0\) . If \(k\) is of characteristic 0, the map \(\lambda\) is linear. If \(k\) is of characteristic \(p \neq 0\) , there exists a power \(q\) of \(p\) dividing \(\dim V^{\lambda}(S)\) , and a homogeneous polynomial function \(P : S \to k\) of degree \(q\) , such that \(\lambda(s)^{q} = P(s)\) for all \(s \in S\) .

Since \(V^{\lambda}(S)\) is stable under \(r(S)\) (Lemma 1 and formula (1) of no. 1), we can assume that \(V = V^{\lambda}(S)\) . Let \(n = \dim V\) . Thus, for \(s \in S\) ,

\[
\det (X - r (s)) = (X - \lambda (s)) ^ {n}.
\]

On the other hand, the expansion of the determinant shows that

\[
\det (\mathrm{X} - r (s)) = \mathrm{X} ^ {n} + a _ {1} (s) \mathrm{X} ^ {n - 1} + \dots + a _ {i} (s) \mathrm{X} ^ {n - i} + \dots
\]

where \(a_{i}: S \to k\) is a homogeneous polynomial function of degree i. Write n = qm where q is a power of the characteristic exponent of k and \((q, m) = 1\) . Then \((\mathrm{X} - \lambda(s))^{n} = (\mathrm{X}^{q} - \lambda(s)^{q})^{m}\) ; hence \(-m\lambda(s)^{q} = a_{q}(s)\) , which implies the result.

PROPOSITION 6. Assume that \(k\) is infinite and that condition (AC) is satisfied. Let \(k'\) be an extension of \(k\) . Put \(V' = V \otimes_k k'\) , \(S' = S \otimes_k k'\) . Let \(r': S' \to \text{End}(V')\) be the map obtained from \(r\) by extension of scalars. Then

\[
\mathrm{V} ^ {0} (\mathrm{S}) \otimes_ {k} k ^ {\prime} = \mathrm{V} ^ {\prime 0} (\mathrm{S}) = \mathrm{V} ^ {\prime 0} (\mathrm{S} ^ {\prime}).
\]

The first equality follows from Prop. 1. To prove the second, we can assume that \(\mathrm{V} = \mathrm{V}^0 (\mathrm{S})\) and so \(\mathrm{V}' = \mathrm{V}'^{0}(\mathrm{S})\) . Let \((s_1,\dots ,s_m)\) be a basis of \(\mathrm{S}\) and \((e_1,\dots ,e_n)\) a basis of \(\mathrm{V}\) . There exist polynomials \(\mathrm{P}_{ij}(\mathrm{X}_1,\dots ,\mathrm{X}_m)\) such that

\[
r ^ {\prime} (a _ {1} s _ {1} + \dots + a _ {m} s _ {m}) ^ {n} e _ {j} = \sum_ {i = 1} ^ {n} \mathrm{P} _ {i j} (a _ {1}, \ldots , a _ {m}) e _ {i}
\]

for \(1 \leq j \leq n\) and \(a_{1}, \ldots, a_{m} \in k'\) . By hypothesis, \(r'(s)^{n} = 0\) for all \(s \in S\) , in other words \(\mathrm{P}_{ij}(a_{1}, \ldots, a_{m}) = 0\) for \(1 \leq i, j \leq n\) and \(a_{1}, \ldots, a_{m} \in k\) . Since k is infinite, \(P_{ij} = 0\) . Consequently, every element of \(r'(S')\) is nilpotent and \(V' = V'^{0}(S')\) .

PROPOSITION 7. Assume that \(k\) is infinite and that condition (AC) is satisfied. Let \(\tilde{\mathrm{S}}\) be the set of \(s \in \mathrm{S}\) such that \(\mathrm{V}^0(s) = \mathrm{V}^0(\mathrm{S})\) . If \(s \in \mathrm{S}\) , let \(\mathrm{P}(s)\) be the determinant of the endomorphism of \(\mathrm{V} / \mathrm{V}^0(\mathrm{S})\) defined by \(r(s)\) (no. 1, Cor. 2 (i) of Th. 1).

\begin{enumerate}
\def\labelenumi{(\roman{enumi})}
\item
  The function \(s \mapsto \mathrm{P}(s)\) is polynomial on \(S\) . We have \(\tilde{\mathrm{S}} = \{s \in S | \mathrm{P}(s) \neq 0\}\) ; this is an open subset of \(S\) in the Zariski topology (App. 1).
\item
  \(\tilde{\mathrm{S}}\) is non-empty, and \(\mathrm{V}^{+}(s) = \mathrm{V}^{+}(\mathrm{S})\) for all \(s\in \tilde{\mathrm{S}}\) .
\end{enumerate}

The fact that \(s \mapsto \mathrm{P}(s)\) is polynomial follows from the linearity of r. If \(s \in \mathrm{S}\) , \(\mathrm{V}^{0}(s) \supset \mathrm{V}^{0}(\mathrm{S})\) , with equality if and only if \(r(s)\) defines an automorphism of \(\mathrm{V}/\mathrm{V}^{0}(\mathrm{S})\) , hence (i).

Now let \(k'\) be an algebraic closure of k, and introduce \(V', S', r'\) as in Prop. 6. We remark that \(S'\) satisfies condition (AC) by continuation of the polynomial identity \((\operatorname{ad} r(s_1))^{2\dim V-1}r(s_2) = 0\) valid for \(s_1, s_2 \in S\) (no. 1, Remark). Applying Th.1, we deduce a decomposition

\[
\mathrm{V} ^ {\prime} = \mathrm{V} ^ {0} (\mathrm{S} ^ {\prime}) \oplus \sum_ {i = 1} ^ {m} \mathrm{V} ^ {\lambda_ {i}} (\mathrm{S} ^ {\prime})
\]

with \(\lambda_{i} \neq 0\) for \(1 \leq i \leq m\) . For \(1 \leq i \leq m\) , there exists a polynomial function \(\mathrm{P}_{i}\) non-zero on \(\mathrm{S}'\) and an integer \(q_{i}\) such that \(\lambda_{i}^{q_{i}} = \mathrm{P}_{i}\) (Prop. 5). Since \(k\) is infinite, there exists \(s \in \mathrm{S}\) such that \((\mathrm{P}_{1} \ldots \mathrm{P}_{m})(s) \neq 0\) , cf.~Algebra, Chap. IV, §2, no. 3, Cor. 2 of Prop. 9. Then \(\lambda_{i}(s) \neq 0\) for all \(i\) , so \(\mathrm{V}^{\prime 0}(\mathrm{S}') = \mathrm{V}^{\prime 0}(s)\) and consequently \(\mathrm{V}^{0}(\mathrm{S}) = \mathrm{V}^{0}(s)\) (Prop. 6), which shows that \(\tilde{\mathrm{S}} \neq \emptyset\) . If \(s \in \tilde{\mathrm{S}}\) , the fact that \(\mathrm{V}^{+}(\mathrm{S})\) is stable under \(r(s)\) and is a complement of \(\mathrm{V}^{0}(s)\) in V implies that \(\mathrm{V}^{+}(\mathrm{S}) = \mathrm{V}^{+}(s)\) (Cor. 2 of Th. 1).

